% fig_spt.tex -- From the physical symmetry assumption to the SPT cohomology invariant.
% Requires: \usepackage{tikz}; \usetikzlibrary{arrows.meta, positioning, calc}.
% Top tier: the physical-layer assumption -- applying the on-site symmetry
%   U(g) to the physical leg of A (defining \tilde A_g) is assumed to generate
%   the same MPV family as A itself ("Sym").
% Middle tier: FT-MPS turns that physical-layer assumption into the
%   virtual-layer gauge identity \tilde A_g^i = X(g) A^i X(g)^{-1}.
%   Tensor-network conventions follow fig_blocking (small filled circles for
%   tensors, thin lines for bonds and physical legs).
% Bottom tier: the chain turning that gauge into a topological invariant.
\begin{figure}[t]
    \centering
    \begin{tikzpicture}[
            font=\scriptsize,
            x=1cm, y=1cm,
            >={Stealth[length=3.2pt]},
            tn dot/.style  ={draw, fill, shape=circle, inner sep=0.055cm},
            tn bond/.style ={line width=0.5pt},
            tn phys/.style ={line width=0.7pt},
            gbox/.style    ={draw, fill=white, rounded corners=1pt, inner sep=1.6pt},
            arr/.style     ={->, semithick, draw=black!80},
            albl/.style    ={font=\scriptsize\itshape, inner sep=1.5pt}
        ]
        %% ---- top tier: physical-layer assumption -----------------------------
        % LHS: tensor A with U(g) on the physical leg, defining \tilde A_g
        \node[tn dot] (P1) at (0,2.62) {};
        \node at (0,2.32) {$A$};
        \draw[tn bond] (P1) -- ++(-0.65,0);
        \draw[tn bond] (P1) -- ++(0.65,0);
        \draw[tn phys] (0,2.68) -- (0,2.88);
        \node[gbox] (Pg) at (0,3.05) {$U(g)$};
        \draw[tn phys] (Pg.north) -- ++(0,0.14);

        \node[align=center] at (1.55,2.62) {\normalsize$\overset{\text{Sym}}{=}$\\[1pt]\scriptsize same MPV};

        % RHS: plain tensor A, for comparison
        \node[tn dot] (P2) at (3.1,2.62) {};
        \node at (3.1,2.32) {$A$};
        \draw[tn bond] (P2) -- ++(-0.65,0);
        \draw[tn bond] (P2) -- ++(0.65,0);
        \draw[tn phys] (3.1,2.68) -- (3.1,3.34);

        % labels sit beside the arrows (not on them, so no white box occludes
        % the shaft or arrowhead)
        \draw[arr] (1.85,2.04) -- node[albl, right=1.5pt]{FT-MPS} (1.85,1.62);

        %% ---- middle tier: virtual gauge identity  \tilde A_g^i = X(g) A^i X(g)^{-1} ----
        \node[tn dot] (L) at (0,1.08) {};
        \node at (0,0.78) {$\tilde A_g$};
        \draw[tn bond] (L) -- ++(-0.65,0);
        \draw[tn bond] (L) -- ++(0.65,0);
        \draw[tn phys] (0,1.14) -- (0,1.42);

        \node at (1.2,1.08) {$=$};

        \node[tn dot] (R) at (3,1.08) {};
        \node at (3,0.78) {$A$};
        \draw[tn phys] (3,1.14) -- (3,1.42);
        \draw[tn bond] (R) -- ++(-0.38,0);
        \node[gbox] (Xg) at (2.30,1.08) {$X(g)$};
        \draw[tn bond] (Xg.west) -- ++(-0.30,0);
        \draw[tn bond] (R) -- ++(0.38,0);
        \node[gbox] (Xgi) at (3.78,1.08) {$X(g)^{-1}$};
        \draw[tn bond] (Xgi.east) -- ++(0.30,0);

        %% ---- bottom tier: gauge -> projective rep -> cohomology class ---------
        % single text lines (no boxes), arrows short with side labels: the same
        % content at less than half the former height
        \draw[arr] (1.85,0.55) -- node[albl, right=1.5pt]{gauge unique up to a scalar} (1.85,0.19);
        \node (proj) at (1.85,-0.05)
            {$\rho(g)\,\rho(h) = \omega(g,h)\,\rho(gh)$: projective rep.\ on the bond space};
        \draw[arr] (1.85,-0.30) -- node[albl, right=1.5pt]{associativity, gauge independence} (1.85,-0.64);
        \node (coh) at (1.85,-0.91)
            {$[\omega]\in H^2(G,\C^\times)$: the topological invariant, the SPT phase label};
    \end{tikzpicture}
    \caption{From the fundamental theorem to the classification of symmetry protected topological phases.
        \emph{Top}: the on-site symmetry $U(g)$ acting on the physical leg of
        an injective MPS tensor $A$ defines $\tilde A_g$; the physical-layer
        assumption is that $\tilde A_g$ generates the same MPV family as $A$
        itself.
        \emph{Middle}: FT-MPS turns this assumption into the virtual-layer
        gauge identity $\tilde A_g^i = X(g)\,A^i\,{X(g)}^{-1}$
        (\cref{thm:spt-invariant}; physical legs drawn upward, bond legs
        horizontal, as in \cref{fig:blocking}).
        \emph{Bottom}: $X(g)$ is unique up to a scalar, so the gauges form a
        projective representation $\rho$ with 2-cocycle $\omega$; its class
        $[\omega]\in H^2(G,\C^\times)$ is the gauge invariant labeling the
        symmetry-protected topological (SPT) phase, $U(1)$-valued after the
        unitary normalization.
        }\label{fig:spt}
\end{figure}

